\documentclass[10pt,a4paper]{amsart}
\usepackage[margin=19mm]{geometry}
\usepackage{amsmath,amssymb,mathtools}
\usepackage[scr=rsfs,cal=euler]{mathalfa}
\usepackage{xfrac}
\usepackage[unicode,hidelinks,bookmarks=false]{hyperref}
\newcommand{\ie}{i.e.\ }
\newcommand{\cf}{cf.\ }
\newcommand{\resp}{respectively\ }
\newcommand{\Subsection}{\subsection}
\providecommand{\nobreakdash}{\nobreak\hbox{-}}
\setlength{\emergencystretch}{2em}
\multlinegap0pt
\usepackage{setspace}
\usepackage{braket}
\usepackage{latexsym}
\usepackage{textcomp}
\usepackage{mathtools}
\usepackage{here}
\newtheorem{Thm}{Theorem}[section]
\newtheorem{Que}[Thm]{Question}
\newtheorem{Lem}[Thm]{Lemma}
\newtheorem{Prop}[Thm]{Proposition}
\newtheorem{Conj}[Thm]{Conjecture}
\newtheorem{Cor}[Thm]{Corollary}
\newtheorem{Cla}[Thm]{Claim}
\newtheorem{Fac}[Thm]{Fact}
\newtheorem{Axi}[Thm]{Axiom}
\newtheorem{Ass}[Thm]{Assumption}
\newtheorem{Def}[Thm]{Definition}
\newtheorem{Rem}[Thm]{Remark}
\newtheorem{Exa}[Thm]{Example}
\begin{document}
\title[The Moran--Hutchinson formula in semimetric spaces]{The Moran--Hutchinson formula in semimetric spaces}
\author{Kazuki Okamura}
\date{}
\maketitle
\noindent\emph{Source and licence.} The Moran--Hutchinson formula in semimetric spaces. Version arXiv:2608.16817v1 (CC BY 4.0). This material is distributed under the Creative Commons Attribution 4.0 International licence, http://creativecommons.org/licenses/by/4.0/, as recorded in the arXiv OAI licence metadata for this exact version and shown on the abstract page https://arxiv.org/abs/2608.16817v1. Full rights notice: licenses/arxiv-2608.16817-CC-BY-4.0.txt.\par\smallskip
\noindent\textit{Editorial scope.} Counted unit: the original \texttt{Thm} environment \texttt{thm:main} at source lines 597--602 together with its complete proof at lines 604--629; the delivered document keeps source lines 137--675 so that every referenced label and every citation resolves inside it. Native editable source is the paper's own concatenated TeX; the AMS wrapper is editorial formatting only. No publisher class, upstream script or unrelated artwork is redistributed.
\begin{Lem}\label{lem:interior-open}
Assume that a semimetric space $(X,d)$ is regular. Then,\\
(i) For every $A \subset X$, $\textup{int}(A)$ is an open subset of $X$. \\
(ii) For every $r > 0$, there exists $\delta > 0$ such that $B(x,\delta) \subset \textup{int}(B(x,r))$ for every $x \in X$. 
\end{Lem}

Part (ii) is the same as \cite[Lemma 2]{KocsisPales2022}. 

\begin{proof}
(i) Let $x \in \textup{int}(A)$. 
Then there exists $\epsilon > 0$ such that $B(x,\epsilon) \subset A$. 
By the regularity, there exists $\delta > 0$ such that $\Phi_d (\delta, \delta) < \epsilon$. 
It suffices to show that $B(x,\delta) \subset  \textup{int}(A)$.  
Let $y \in B(x,\delta)$. 
Then, for every $z \in B(y,\delta)$, 
\[ d(x,z) \le \Phi_d (d(x,y), d(y,z)) \le \Phi_d (\delta, \delta) < \epsilon. \]
Hence, $z \in B(x,\epsilon) \subset A$. 
Hence $B(y,\delta) \subset A$, which means $y \in \textup{int}(A)$. 

Part (ii) follows by the same argument. 
\end{proof}

By Lemma \ref{lem:interior-open}, we can show that every regular semimetric space is Hausdorff. 

The following result means that the topological and sequential closures of a set coincide. 

\begin{Lem}\label{lem:seq-closure}
If a semimetric space $(X,d)$ is regular, then for every $H \subset X$, 
\begin{equation}\label{eq:closure-limit} 
\overline{H} = \left\{x \in X \colon \textup{there exists a sequence $(x_n)_n$ in $H$ such that $\lim_{n \to \infty} d(x_n, x) = 0$} \right\}. 
\end{equation}
\end{Lem}

\begin{proof}
Let $\widetilde{H}$ be the set on the right-hand side of \eqref{eq:closure-limit}. 
$\widetilde{H}$ contains $H$. 
We show that $\widetilde{H}$ is closed. 
This is equivalent to saying that $X \setminus \widetilde{H}$ is open. 
Let $x \in X \setminus \widetilde{H}$. 
Assume that for every $n \ge 1$, there exists $x_n \in \widetilde{H} \cap B(x, 2^{-n})$. 
Then, for every $n \ge 1$, there exists $y_n \in H$ such that $d(x_n, y_{n}) \le 2^{-n}$. 
By the regularity, 
\[ d(y_{n}, x) \le \Phi_d \left(d(y_{n}, x_n), d(x_n, x)\right) \le \Phi_d (2^{-n}, 2^{-n}) \to 0, \quad n \to \infty. \]
Hence $x \in \widetilde{H}$. 
This is a contradiction. 
Hence there exists $n \ge 1$ such that $ \widetilde{H} \cap B(x, 2^{-n}) = \emptyset$. 
Thus we see that $X \setminus \widetilde{H}$ is open and that $\overline{H} \subset \widetilde{H}$. 

Let $x \in \widetilde{H}$. 
Let $U$ be an open subset of $X$ containing $x$. 
Then there exists $r > 0$ such that $B(x,r) \subset U$ and there exists $y \in H$ such that $y \in B(x,r)$. 
Hence $H \cap U \ne \emptyset$. 
Hence $x \in \overline{H}$. 
Thus we see that $\widetilde{H} \subset \overline{H}$. 
\end{proof}

By this lemma, a subset of a complete regular semimetric space is complete if and only if it is closed. 
A regular semimetric space  $(X,d)$ is compact if and only if it is sequentially compact. 


We now state some consequences of the normal property. 
We say that a subset $H$ of $X$ is {\it bounded} if there exist $x_0 \in X$ and $r_0 > 0$ such that $H \subset B(x_0, r_0)$. 
Let 
\[ \textup{diam}(H) \coloneqq \sup\left\{d(x,y) | x, y \in H \right\} \] 
for $H \subset X$, $H \ne \emptyset$, and we call it the {\it diameter} of $H$. 

\begin{Lem}[{\cite[Proposition 3]{BessenyeiPenzes2022}}]\label{lem:bounded-basic}
Let $(X,d)$ be a normal semimetric space and $H$ be a nonempty subset of $X$. 
Then, $H$ is bounded if and only if $\textup{diam}(H) < \infty$. 
\end{Lem}

We say that a subset $H$ of $X$ is {\it totally bounded} if for every $\epsilon > 0$ there exist finitely many points $x_1, \dots, x_n \in H$ such that $H \subset \bigcup_{i=1}^{n} B(x_i, \epsilon)$. 
By \cite[Lemma 2]{BessenyeiPenzes2022}, 
we see that if $(X,d)$ is complete and regular and $H \subset X$, then $H$ is compact if and only if $H$ is closed and totally bounded. 

{\it Hereafter, we assume that $(X,d)$ is a complete semimetric space satisfying the regular and normal properties for the basic triangle function.}  



We adopt Matkowski-Rus's definition of weak contractions \cite{Matkowski1975,Rus2001}.  
We say that a map $\phi : [0,\infty) \to [0,\infty)$ is a {\it comparison function} if 
$\phi$ is nondecreasing, and 
$\lim_{n \to \infty} \phi^n (t) = 0$ for every $t > 0$. 
Let $\phi$ be a comparison function. 
We say that a map $T : X \to X$ is a {\it $\phi$-contraction} if 
$d(Tx, Ty) \le \phi(d(x,y))$ for every $x, y \in X$. 



We have the following. 
\begin{Lem}\label{lem:comparison-basic}
(i) If $\phi$ is a comparison function, then, $\phi(0) = 0$ and  $\phi(t) < t$ for every $t > 0$. \\
(ii) If $\phi_1$ and $\phi_2$ are right-continuous comparison functions and $\phi \coloneqq \max\{\phi_1, \phi_2\}$, then, $\phi$ is also a right-continuous comparison function. 
\end{Lem}

We now recall some consequences of the regular property. 

\begin{Thm}[Bessenyei and P\'ales {\cite[Theorem 1]{BessenyeiPales2017}} ]\label{thm:Bessenyei-Pales-FP} 
Let $\phi$ be a comparison function and let $T : X \to X$ be a $\phi$-contraction. 
Then, there exists a unique fixed point of $T$. 
\end{Thm}

Let $\mathcal{K}(X)$ be the set of nonempty compact subsets of $X$. 

\begin{Thm}[Kocsis and P\'ales {\cite[Theorem 17]{KocsisPales2022}}; Bessenyei and P\'{e}nzes {\cite[Theorem 2]{BessenyeiPenzes2022}}]\label{thm:attractor}
Let $\ell \ge 1$. 
For each $i \in \{1, \dots, \ell\}$, 
let $f_i$ be a $\phi_i$-contraction on $X$ where $\phi_i$ is a right-continuous comparison function.
Then, there exists a unique $K \in \mathcal{K}(X)$ such that $K = \bigcup_{i=1}^{\ell} f_i (K)$. 
\end{Thm}

We call the set $K$ in the above theorem the {\it attractor} of $\{f_i\}_i$. 


\begin{Rem}
A semimetric space $(X,d)$ is called a {\it quasi-metric space} if there exists $C < \infty$ such that $d(x,z) \le C(d(x,y) + d(y,z))$ for every $x, y, z \in X$. 
This is also known as a {\it $b$-metric space} \cite{CJT2018}. 
Every quasi-metric space is normal and regular. 
Properties of quasi-metric spaces were investigated by Bourbaki \cite{Bourbaki2007}. 
\end{Rem}




We follow \cite{BessenyeiPenzes2022} for the definitions of Hausdorff measure and Hausdorff dimension. 
These are defined in the same manner as in the case of metric spaces. 

For $A \subset X$, $\alpha \ge 0$ and $\delta > 0$, let 
\[ \mathcal{H}^{\alpha}_{\delta}(A) \coloneqq \inf\left\{ \sum_{i=1}^{\infty} \textup{diam}(B_i)^{\alpha} \colon A \subset \bigcup_{i} B_{i}, \textup{diam}(B_i) < \delta   \right\} \]
and $\mathcal{H}^{\alpha}(A) \coloneqq \lim_{\delta \to +0} \mathcal{H}^{\alpha}_{\delta}(A)$. 
The Hausdorff dimension of $A \subset X$ is defined by 
\[ \dim_H (A) \coloneqq \sup\{r > 0 \colon \mathcal{H}^r (A) = \infty \}. \]
We see that $\mathcal{H}^t (A) = 0$ for every $t > r$ whenever $\mathcal{H}^r (A) < \infty$ and that $\dim_H (A) = \inf\{r > 0 \colon \mathcal{H}^{r}(A) = 0\}$. 

We define the Minkowski dimension in the same manner as in the case of metric spaces. 
For a totally bounded subset $A$ of $X$, 
let $N(A, \epsilon)$ be the minimum number of sets of diameter less than $\epsilon$ needed to cover $A$. 
Define the upper and lower Minkowski dimensions by 
\[ \overline{\dim}_M A \coloneqq \limsup_{\epsilon \to +0} \frac{\log N(A, \epsilon)}{\log 1/\epsilon},  \textup{ and } \underline{\dim}_M A \coloneqq \liminf_{\epsilon \to +0} \frac{\log N(A, \epsilon)}{\log 1/\epsilon},  \]
respectively. 
If these two values agree, then we call their common value the {\it Minkowski dimension} of $A$ and denote it by $\dim_M A$. 
Then, 
\[ \dim_H A \le \underline{\dim}_M A \le \overline{\dim}_M  A.  \]



\section{Existence of an invariant measure}\label{sec:inv-measure}

The goal of this section is to show the following. 

\begin{Prop}\label{prop:exist-inv-measure}
Assume that $p_1, \dots, p_{\ell} > 0$, $\sum_{j = 1}^{\ell} p_j = 1$. 
For each $j \in \{1, \dots, \ell\}$, 
let $f_j$ be a $\phi_j$-contraction on $X$ where $\phi_j$ is a right-continuous comparison function. 
Let $K$ be the attractor of $\{f_j\}_j$. 
Then, there exists a Borel probability measure $\mu$ supported on $K$ such that 
\begin{equation}\label{eq:invariance-measure} 
\mu = \sum_{j = 1}^{\ell} p_j \mu \circ f_j^{-1}. 
\end{equation}
\end{Prop}

For $\xi = (\xi_i)_i \in \{1, \dots, \ell\}^{\mathbb N}$, 
let  $f_{\xi|_{n}} \coloneqq f_{\xi_1} \circ \cdots \circ f_{\xi_n}$. 

\begin{Lem}
The set $\bigcap_{n \ge 1} f_{\xi|_{n}}(K)$ is a singleton contained in $K$. 
\end{Lem}

\begin{proof}
Since each comparison function $\phi_j$ is right-continuous, each $f_j$ is continuous. 
Hence $f_{\xi|_{n}}$ is continuous. 
Since $K$ is compact, $f_{\xi|_{n}}(K)$ is also compact. 
Since $f_j (K) \subset K$ for each $j$, the sets $f_{\xi|_{n}}(K)$ form a decreasing sequence. 
Hence, $\bigcap_{n \ge 1}  f_{\xi|_{n}}(K)$ is not empty. 

Let $\phi \coloneqq \max_{j \in \{1, \dots, \ell\}} \phi_j$. 
Then, by Lemma \ref{lem:comparison-basic}, this is also a right-continuous comparison function. 
We see that by induction on $n$, 
$\textup{diam}( f_{\xi|_{n}}(K)) \le \phi^n (\textup{diam}(K))$. 
Since $K$ is compact, by Lemma \ref{lem:bounded-basic}, $\textup{diam}(K) < \infty$.  
 
Let $x_1, x_2 \in \bigcap_{n \ge 1}  f_{\xi|_{n}}(K)$. 
Then, $d(x_1, x_2) \le \textup{diam}( f_{\xi|_{n}}(K))$ for each $n$. 
Hence, $d(x_1, x_2) = 0$ and then $x_1 = x_2$. 
\end{proof}

For ease of notation, let $\Omega \coloneqq \{1, \dots, \ell\}^{\mathbb N}$. 
Define a map $\Psi \colon \Omega \to K$ by $\{\Psi(\xi)\} = \bigcap_{n \ge 1} f_{\xi|_{n}}(K)$.  
We equip $\Omega$ with the product topology. 

\begin{Lem}\label{lem:conti-surj}
The map $\Psi$ is a continuous surjection. 
\end{Lem}

\begin{proof}
First we show that $\Psi$ is continuous. 
For $\xi = (\xi_i)_{i \ge 1}, \eta = (\eta_i)_{i \ge 1} \in \Omega$, 
let $|\xi \wedge \eta| \coloneqq \min\{i \ge 1 \colon \xi_i \ne \eta_i\}$ and $d_{\Omega} (\xi, \eta) \coloneqq \exp(-|\xi \wedge \eta|)$. 
We remark that $|\xi \wedge \eta| = \infty$ if and only if $\xi = \eta$. 
We assume that $\min \emptyset = \infty$ and $\exp(-\infty) = 0$. 
This defines a metric on $\Omega$ and the metric topology is identical to the product topology. 
For $\xi \ne \eta$, we see that 
\[ d(\Psi(\xi), \Psi(\eta)) \le \textup{diam}\left(f_{\xi|_{n-1}} (K) \right) \le \phi^{n -1}(\textup{diam}(K)) \]
\[ = \phi^{-\log d_{\Omega}(\xi, \eta) -1}(\textup{diam}(K)), \quad n = |\xi \wedge \eta|,  \]
where we let $\xi|_0 \coloneqq \emptyset$ and recall that $f_{\emptyset}$ is the identity map on $X$. 
Since $\phi$ is a comparison function, 
for every $\epsilon > 0$, there exists $\delta > 0$ such that $d_{\Omega}(\xi, \eta) < \delta$ implies $d(\Psi(\xi), \Psi(\eta)) < \epsilon$. 
Hence $\Psi$ is continuous. 

We now show that $\Psi$ is surjective. 
Let $x \in K$. 
We define a sequence $\xi = (\xi_i)_i \in \Omega$ inductively. 
Since $K$ is the attractor, there exists $j  \in \{1, \dots, \ell\}$ such that $x \in f_{j} (K)$, and we let $\xi_1 \coloneqq j$ by choosing one such $j$. 
We choose $x_1 \in K$ such that $x = f_{\xi_1} (x_1)$. 
In the same manner, there exists $j \in \{1, \dots, \ell\}$ such that $x_1 \in f_{j} (K)$, and we let $\xi_2 \coloneqq j$ by choosing one such $j$. 
We choose $x_2 \in K$ such that $x_1 = f_{\xi_2} (x_2)$. 
We repeat this procedure and obtain a sequence $\xi = (\xi_i)_i \in \Omega$. 
 Then, for every $n \ge 1$, $x \in f_{\xi|_{n}} (K)$ and hence $x = \Psi(\xi)$. 
\end{proof}

Let $\iota_{\bf p}$ be the probability measure on $\{1,\dots, \ell\}$ such that $\iota_{\bf p} (\{j\}) = p_j$ for each $j$. 
Let $\nu_{\bf p}$ be the product measure $\iota_{\bf p}^{\otimes \mathbb N}$ on $\Omega$. 
By Lemma \ref{lem:conti-surj}, $\Psi$ is a Borel measurable map. 
Let $\mu_{\bf p}$ be the push-forward measure of $\nu_{\bf p}$ by the map $\Psi$.  
This is a Borel probability measure on $K$. 
We regard it as a probability measure on $X$, extended by zero outside $K$. 

We show that $\mu_{\bf p} = \sum_{j = 1}^{\ell} p_j \mu_{\bf p} \circ f_{j}^{-1}$. 
Let $A$ be a Borel measurable subset of $K$. 
Let $\tilde{\xi} \coloneqq (\xi_{i+1})_{i \ge 1}$ for $\xi = (\xi_i)_{i \ge 1}$.  
Then, $\Psi(\xi) = f_{\xi_1}\left(\Psi (\tilde{\xi}) \right)$.  
Therefore, 
\[ \mu_{\bf p} (A) = \nu_{\bf p}  (\Psi^{-1}(A)) = \sum_{j = 1}^{\ell} \nu_{\bf p}  \left(  \left\{\xi \in \{1, \dots, \ell\}^{\mathbb N}  \colon \xi_1 = j, \, \Psi(\xi) \in A \right\} \right)  \] 
\[ = \sum_{j = 1}^{\ell} \nu_{\bf p}  \left( \left\{\xi  \in \{1, \dots, \ell\}^{\mathbb N} \colon \xi_1 = j, \, \tilde{\xi} \in \Psi^{-1}(f_j^{-1}(A)) \right\}  \right) \]
\[ = \sum_{j = 1}^{\ell} p_j \nu_{\bf p}  \left(\Psi^{-1}(f_j^{-1}(A))  \right) =  \sum_{i=1}^{\ell} p_i \mu_{\bf p} (f_i^{-1} (A)).  \]
Since $f_j (K) \subset K$ for each $j$, $\mu_{\bf p} \circ f_j^{-1}$ vanishes on $K^c$. 
Hence \eqref{eq:invariance-measure} holds for every Borel measurable subset of $X$. 

\begin{Rem}
We have {\it not} shown the uniqueness of the invariant measure satisfying \eqref{eq:invariance-measure}. 
\end{Rem}



\section{Moran--Hutchinson formula}\label{sec:MH-formula}

In this section we establish the Moran--Hutchinson formula for the attractor of a finite family of similitudes. 
As an outline, we follow the strategy in the proof of  \cite[Theorem 2.2.2]{BishopPeres2017}. 
However, we need to modify several parts of the proof.  


\begin{Ass}\label{ass:two-additional-metric}
(i) (strong regularity)
\[ \limsup_{\epsilon \to +0} \frac{\Phi_d (\epsilon, \epsilon)}{\epsilon} < \infty.  \]
(ii) (geometric doubling)
There exists $N_0 \in \mathbb{N}$ such that for every $x \in X$ and every $r > 0$, there exist $x_1, \dots, x_{N_0} \in X$ such that $B(x,r) \subset \bigcup_{i=1}^{N_0} B(x_i, r/2)$. 
\end{Ass}

Every quasi-metric space satisfies the strong regularity above. 
The following corresponds to \cite[Lemma 3.3]{BaloghRohner2007}. 

\begin{Lem}\label{lem:number-disjoint}
Suppose Assumption \ref{ass:two-additional-metric} holds. 
Let $b > a > 0$. 
Let $r > 0$. 
Let $W_1, \dots, W_N$ be disjoint subsets of $X$ such that for each $i$, $W_i$ contains a ball of radius $ar$ and is contained in a ball of radius $br$. 
Assume that there exists a subset $A$ of $X$ such that  
$\textup{diam}(A) \le r$ and $A \cap W_i \ne \emptyset$ for each $i \in \{1, \dots, N\}$. 
Then, there exist  $n_0 (a,b) \in \mathbb{N}$ and $R_0 (a,b) > 0$ such that whenever $r \in (0, R_0 (a,b))$ and the above conditions hold, $N \le N_0^{n_0 (a,b)}$. 
\end{Lem}

\begin{proof}
By the assumption, for each $i$, there exist two points $y_i, z_i \in X$ such that $B(y_i, ar) \subset W_i \subset B(z_i, br)$. 

Step 1. We show that there exist $x \in A$  and  two constants $C(b) >  1$ and $R_1 (b) > 0$ such that $W_i \subset B(x, C(b) r)$ for each $i$ and every $r \in (0, R_0 (b))$. 

Take a point $x \in A$. 
Let $x_i \in A \cap W_i$.   
Then, for every $w \in W_i$, 
\[ d(x,w) \le \Phi_d (d(x,x_i), d(x_i,w)) \le \Phi_d \left(r, \Phi_d (br, br)\right)\]
\[  \le  \Phi_d \left(u,  u\right), \quad u = \max\left\{r, \Phi_d (br,br)\right\}. \]
By Assumption \ref{ass:two-additional-metric} (i), 
there exist $R_1 > 0$ and $C_1 > 1$ such that for every $r \in (0, R_1)$, $\Phi_d (r,r) \le C_1 r$. 
The claim follows for $R_1 (b) \coloneqq R_1/(1+ C_1 b)$ and $C(b) \coloneqq C_1 (1+ C_1 b) $. 

Step 2. For every $r \in (0, R_0 (b))$, $B(y_i, ar), 1 \le i \le N$, are disjoint balls  contained in $B(x, C(b) r)$. 
Then there exist $n_0 = n_0 (a,b) \in \mathbb{N}$ and $R_2 (a,b) > 0$ such that for every $r \in (0, R_2 (a,b))$,  $\Phi_d (C(b) r/2^{n_0}, C(b) r/2^{n_0}) < ar$. 

Let $R_0 (a,b) \coloneqq \min\{R_1 (b), R_2 (a,b) \}$. 
Fix $r \in \left(0, R_0 (a,b)\right)$. 
Then by Assumption \ref{ass:two-additional-metric} (ii), 
there exist $v_1, \dots, v_{N_0^{n_0}} \in X$ such that $B(x,C(b)r) \subset \bigcup_{k=1}^{N_0^{n_0}} B(v_k, C(b) r/2^{n_0})$. 
If $N > N_0^{n_0}$, then by the pigeonhole principle, there exist $i \ne j$ and $k$ such that $y_i \in B(v_k, C(b) r/2^{n_0})$ and $y_j \in B(v_k, C(b) r/2^{n_0})$. 
It holds that 
\[  d(y_i, y_j) \le \Phi_d (d(y_i, v_k), d(v_k, y_j)) \le \Phi_d (C(b) r/2^{n_0}, C(b) r/2^{n_0}) < ar. \]
This contradicts the assumption that $B(y_i, ar), 1 \le i \le N$, are disjoint. 
Hence $N \le N_0^{n_0}$. 
\end{proof}

\begin{Rem}
In the above proof, the geometric doubling condition in Assumption \ref{ass:two-additional-metric} (ii) can be weakened to a {\it local} geometric doubling condition: 
there exists $N_0 \in \mathbb{N}$ and $r_0 > 0$ such that for every $x \in X$ and every $r \in (0, r_0)$, there exist $x_1, \dots, x_{N_0} \in X$ such that $B(x,r) \subset \bigcup_{i=1}^{N_0} B(x_i, r/2)$. 
\end{Rem}



\begin{Ass}\label{ass:additional-f}
(i) (open set condition)  There exists a nonempty, bounded open subset $V$ of $X$ such that $\bigcup_j f_j (V) \subset V$ and 
$f_i (V) \cap f_j (V) = \emptyset$ for $i \ne j$. \\
(ii) (surjectivity and similarity) For each $j$, $f_j$ is a surjective map and a similitude with contraction ratio $r_j$, specifically, 
there exists $r_j \in (0,1)$ such that $d(f_j (x_1), f_j (x_2)) = r_j d(x_1, x_2)$ for every $x_1, x_2 \in X$. 
\end{Ass}

If $(X,d)$ is a Euclidean space and $f$ is a similitude on $X = \mathbb{R}^d$ such that $d(f(x_1), f(x_2)) = r d(x_1, x_2)$ for every $x_1, x_2 \in X$, 
then there exists a $d \times d$ orthogonal matrix $O$ and $x_0 \in X$ such that $f(x) = rOx + x_0$ for every $x \in X$. 
In particular $f$ is bijective. 

\begin{Lem}\label{lem:K-included}
Under Assumption \ref{ass:additional-f} (i), $K \subset \overline{V}$. 
\end{Lem}

\begin{proof}
Since $f_j$ is continuous, $f_j (\overline{V}) \subset \overline{f_j (V)}$. 
Hence  
\[ \bigcup_j f_j (\overline{V}) \subset \bigcup_j \overline{f_j (V)} \subset \overline{V}.\]  
Since $V$ is not empty, we can take a point $v_0 \in V$. 
Since $K$ is compact, by recalling that $(X,d)$ is normal, $M \coloneqq \sup_{x \in K} d(x, v_0) < \infty$. 
Let $x \in K$. 
Then  there exists $\xi \in \{1,\dots, \ell\}^{\mathbb N}$ such that $\Psi(\xi) = x$. 
For every $n$, $x \in f_{\xi|_n} (K)$. 
Hence, $d(x,  f_{\xi|_n}(v_0)) \le \phi^n (M)$. 
Since $ f_{\xi|_n}(v_0) \in V$ for every $n$, we see that $x \in \overline{V}$. 
By \eqref{eq:closure-limit}, we see that $K \subset \overline{V}$. 
\end{proof}

Let $r_{\min} \coloneqq \min_j r_j \in  (0,1), r_{\max} \coloneqq \max_j r_j \in (0,1)$. 
Let $\{1,\dots, \ell\}^{*} \coloneqq \bigcup_{n \ge 0} \{1,\dots, \ell\}^{n}$, where we let $\{1,\dots,\ell\}^{0} \coloneqq \{\emptyset\}$. 
We let $r_{\emptyset} \coloneqq 1$ and $f_{\emptyset}$ be the identity map of $X$. 
For $\sigma = (\sigma_i)_{i=1}^{m} \in  \{1,\dots, \ell\}^{*}$, 
let $|\sigma| \coloneqq m$, $r_{\sigma} \coloneqq r_{\sigma_1} \cdots r_{\sigma_m}$, $p_{\sigma} \coloneqq p_{\sigma_1} \cdots p_{\sigma_m}$ and $f_{\sigma} \coloneqq f_{\sigma_1} \circ \cdots \circ f_{\sigma_m}$. 

\begin{Lem}\label{lem:ball-surjection}
Suppose Assumption \ref{ass:additional-f} (ii) holds. 
Then, for every $\sigma \in \{1,\dots, \ell\}^*$, every $\delta > 0$ and every $x \in X$, 
$B(f_{\sigma}(x), \delta r_{\sigma}) = f_{\sigma}(B(x,\delta))$. 
\end{Lem}

By this lemma, $f_j$ is an open map and hence $f_j$ is a homeomorphism of $X$ onto itself. 

\begin{proof}
Let $y \in B(f_{\sigma}(x), \delta r_{\sigma})$. 
Since $f_{\sigma}$ is surjective, 
there exists $z$ such that $y = f_{\sigma}(z)$. 
We see that $\delta r_{\sigma} > d(f_{\sigma}(x), f_{\sigma}(z)) = r_{\sigma} d(x,z)$.
Hence $d(x,z) < \delta$. 
Hence $B(f_{\sigma}(x), \delta r_{\sigma}) \subset f_{\sigma}(B(x,\delta))$. 
Let $z \in B(x,\delta)$. 
Then $d(f_{\sigma}(x), f_{\sigma}(z)) = r_{\sigma} d(x,z) < r_{\sigma} \delta$. 
\end{proof}


We say that $\Pi \subset \{1,\dots, \ell\}^{*}$ is a {\it cut-set} if for every $\xi \in \{1,\dots, \ell\}^{\mathbb N}$ there exists $\sigma \in \Pi$ such that $\sigma$ is a prefix of $\xi$. 
We say that a cut-set is {\it minimal} if there exists no element in the cut-set which is a prefix of another element. 
Then, by \eqref{eq:invariance-measure}, 
we obtain that for every minimal cut-set $\Pi$, 
\begin{equation}\label{eq:invariance-min-cutset}
\mu_{\bf p} = \sum_{\sigma \in \Pi} p_{\sigma} \mu_{\bf p} \circ f_{\sigma}^{-1}.
\end{equation}
See \cite[Lemma 2.2.4 (ii)]{BishopPeres2017}. 

For $\rho \in (0, r_{\min})$, 
let 
\begin{equation}\label{eq:def-cutset-1} 
\Pi_{\rho} \coloneqq \left\{\sigma \in \{1,\dots, \ell\}^{*} \colon r_{\sigma} \le \rho < r_{\sigma^{\prime}} \right\}, 
\end{equation}
where $\sigma^{\prime} \coloneqq (\sigma_1, \dots, \sigma_{|\sigma|-1})$. 
Since $r_{\max} < 1$, for every $\xi \in \{1,\dots, \ell\}^{\mathbb N}$, 
$\lim_{n \to \infty} r_{\xi|_n} = 0$. 
Hence there exists a unique $n$ such that $r_{\xi|_n} \le \rho < r_{\xi|_{n-1}}$, and then $\xi|_n \in \Pi_{\rho}$. 
Let $\sigma, \tau \in \Pi_{\rho}$, and suppose that $\sigma$ is a prefix of $\tau$. 
Then $|\sigma| \le |\tau|$. 
If $|\sigma| < |\tau|$, then $r_{\tau^{\prime}} \le r_{\sigma}$. 
This cannot occur. 
Hence, $|\sigma| = |\tau|$ and then $\sigma = \tau$. 
Thus we see that $\Pi_{\rho}$ is a minimal cut-set. 




\begin{Prop}\label{prop:mass-dist}
Suppose that Assumptions \ref{ass:two-additional-metric} and \ref{ass:additional-f}  hold. 
Let $\alpha$ be a positive real number such that $\sum_j r_j^{\alpha} = 1$. 
Let $p_j \coloneqq r_j^{\alpha}$ for each $j$ and ${\bf p} \coloneqq (p_j)_j$. 
Then there exist $C_0 \in (0,\infty)$ and $\delta_0 > 0$ such that 
$\mu_{\bf p}(U) \le C_0 \rho^{\alpha}$ for every $\rho \in (0, \delta_0)$ and every open set $U$ with $\textup{diam}(U) \le \rho$. 
\end{Prop}

\begin{proof}
Let $b \coloneqq 1 + \textup{diam}(V)$. 
Let $\rho \in \left(0, \min\{R_0 (b), r_{\min} \}\right)$. 
By \eqref{eq:invariance-min-cutset}, we obtain that 
\[ \mu_{\bf p}(U) =  \sum_{\sigma \in \Pi_{\rho}} r_{\sigma}^{\alpha} \mu_{\bf p} \left( f_{\sigma}^{-1} (U) \right) \le \rho^{\alpha} \left|\left\{\sigma \in \Pi_{\rho} \colon \mu_{\bf p} \left( f_{\sigma}^{-1} (U) \right) > 0 \right\}\right|.  \]
If $\mu_{\bf p} \left( f_{\sigma}^{-1} (U) \right) > 0$, then $f_{\sigma}^{-1} (U) \cap K \ne \emptyset$. 
By Lemma \ref{lem:K-included}, $f_{\sigma}^{-1} (U) \cap \overline{V} \ne \emptyset$. 
Since $U$ is open and $f_{\sigma}$ is continuous, 
$f_{\sigma}^{-1} (U) \cap V \ne \emptyset$, and hence $U \cap f_{\sigma}(V) \ne \emptyset$. 

Since $\Pi_{\rho}$ is a minimal cut-set, by the open set condition, 
$\{ f_{\sigma}(V) \}_{\sigma \in \Pi_{\rho}}$ are disjoint. 
Since $V$ is a nonempty bounded open set, 
there exist $z \in V$ and $\delta > 0$ such that 
$B(z,\delta) \subset V \subset B\left(z, 1+\textup{diam}(V) \right)$. 
We remark that by the normal property, $\textup{diam}(V) < \infty$.  
By Lemma \ref{lem:ball-surjection}, 
\[ B\left(f_{\sigma}(z), \rho \delta r_{\min}\right) \subset B(f_{\sigma}(z), \delta r_{\sigma}) \subset f_{\sigma}(V) \subset B\left(f_{\sigma}(z), \rho(1+ \textup{diam}(V))\right).\] 
 We apply Lemma \ref{lem:number-disjoint} to the case where $a = \delta r_{\min}$ and $b = 1+\textup{diam}(V)$ and $r = \rho$. 
 Then, $|\{\sigma \in \Pi_{\rho} \colon U \cap f_{\sigma}(V) \ne \emptyset\}| \le N_0^{n_0 (a,b)}$. 
 Now the assertion holds for $C_0 =  N_0^{n_0 (a,b)}$ and $\delta_0 = \min\{R_0 (b), r_{\min} \}$. 
\end{proof}

\begin{Cor}\label{cor:non-atomic}
Under the assumptions of Proposition \ref{prop:mass-dist}, $\mu_{\bf p}$ is non-atomic.
\end{Cor}

\begin{proof}
Let $x \in X$. 
By Assumption \ref{ass:two-additional-metric} (i), 
there exist $C_1 > 0$ and $\delta_1 \in (0, \delta_0)$ such that for every $\delta \in (0, \delta_1)$, 
$\textup{diam}(\textup{int}(B(x,\delta))) \le \Phi_d (\delta, \delta) \le C_1 \delta$. 
By Lemma \ref{lem:interior-open} (i), $\textup{int}(B(x,\delta))$ is open.
By Lemma \ref{lem:interior-open} (ii), there exists $\eta > 0$ such that $B(x,\eta) \subset \textup{int}(B(x,\delta))$. 
Then by Proposition \ref{prop:mass-dist}, 
\[ \mu_{\bf p}(\{x\}) \le \mu_{\bf p} \left( \textup{int}(B(x,\delta)) \right) \le C_0 (C_1 \delta)^{\alpha}.   \]
By letting $\delta \to +0$, we have the assertion. 
\end{proof}

We remark that a ball may not be a Borel set, so $\mu_{\bf p}(B(x,\eta))$ might not be well-defined. 

We need the following lemma since a ball may not be open. 

\begin{Lem}\label{lem:large-open}
There exist $C_2 > 0$ and $\delta_2 \in (0, \delta_1)$  such that for every subset $A$ of $X$ with $0 < \textup{diam}(A) < \delta_2$, 
there exists an open subset $U$ such that $A \subset U$ and $ \textup{diam}(U) \le C_2 \textup{diam}(A)$. 
\end{Lem}

\begin{proof}
Assume that $0 < \textup{diam}(A) < \infty$. 
Since $A \ne \emptyset$, we can take a point $x \in A$. 
Then, by using $\textup{diam}(A) > 0$, we obtain that $A \subset B(x, 2\textup{diam}(A))$. 
Take $C_1 > 0$ and $\delta_1$ as in the proof of Corollary \ref{cor:non-atomic}. 

If $0 < \textup{diam}(A) < \delta_1 / (4(C_1 + 1))$, then, $\Phi_d (2\textup{diam}(A), 2\textup{diam}(A)) \le 2C_1 \textup{diam}(A)$. 
By using the same argument as in the proof of Lemma \ref{lem:interior-open} (i), 
$B(x, 2\textup{diam}(A)) \subset \textup{int} \left(B(x, 4C_1 \textup{diam}(A)) \right)$. 
Then 
\[ \textup{diam}\left(  \textup{int} \left(B(x, 4C_1 \textup{diam}(A)) \right) \right) \le \Phi_d (4C_1 \textup{diam}(A),  4C_1 \textup{diam}(A)) \le 4C_1^2 \textup{diam}(A). \]
Now the assertion holds for $\delta_2 =  \delta_1 / (4(C_1 + 1))$, $C_2 = 4C_1^2$, and $U =  \textup{int} \left(B(x, 4C_1 \textup{diam}(A)) \right)$. 
\end{proof}



\begin{Thm}\label{thm:main}
Suppose that  Assumptions \ref{ass:two-additional-metric} and \ref{ass:additional-f}  hold. 
Let $K$ be the attractor of $\{f_j\}_j$. 
Let $\alpha$ be the similarity dimension of $\{f_j\}_j$, specifically,  a positive real number such that $\sum_j r_j^{\alpha} = 1$. 
Then $0 < \mathcal{H}^{\alpha}(K) < \infty$ and $\dim_H K = \dim_M K = \alpha$. 
\end{Thm}

\begin{proof}
We show that $ \mathcal{H}^{\alpha}(K) > 0$. 
The following arguments correspond to the proof of the mass distribution principle, but we need several modifications. 

Let $\delta \in (0, \delta_2 /(C_2 + 1))$. 
Assume that $K \subset \bigcup_{i \ge 1} A_i$ and that $\textup{diam}(A_i) < \delta$ for every $i$. 
Let $\mathcal{I} \coloneqq \{i \ge 1 \colon \textup{diam}(A_i) = 0\}$.   
By Lemma \ref{lem:large-open}, 
it holds that for every $i \notin \mathcal{I}$, there exists an open set $U_i$ such that $A_i \subset U_i$ and $\textup{diam}(U_i) \le C_2 \textup{diam}(A_i)$. 
Since 
\[ K \setminus \bigcup_{i \in \mathcal{I}} A_i \subset  \bigcup_{i \ge 1, i \notin \mathcal{I}} A_i \subset  \bigcup_{i \ge 1, i \notin \mathcal{I}} U_i, \]
we obtain that 
\[ \mu_{\bf p}\left( K \setminus \bigcup_{i \in \mathcal{I}} A_i  \right) \le \sum_{i \ge 1, i \notin \mathcal{I}} \mu_{\bf p}(U_i). \]

We remark that for every $i \in \mathcal{I}$, $A_i$ is empty or a one-point set. 
By Corollary \ref{cor:non-atomic}, $ \mu_{\bf p}\left( K \setminus \bigcup_{i \in \mathcal{I}} A_i  \right)  =  \mu_{\bf p}\left( K \right)  = 1$. 
Since $ \textup{diam}(U_i) \le C_2 \textup{diam}(A_i) < \delta_2 < \delta_0$, 
we can apply Proposition \ref{prop:mass-dist} and obtain that for $i \notin \mathcal{I}$, 
\[ \mu_{\bf p}(U_i) \le C_0  \textup{diam}(U_i) ^{\alpha} \le C_0 C_2^{\alpha} \textup{diam}(A_i)^{\alpha}.\] 
Then, 
\[ \sum_{i \ge 1} \textup{diam}(A_i)^{\alpha} = \sum_{i \ge 1, i \notin \mathcal{I}} \textup{diam}(A_i)^{\alpha} \ge \frac{1}{C_0 C_2^{\alpha}} \sum_{i \ge 1, i \notin \mathcal{I}} \mu_{\bf p}(U_i) \ge  \frac{1}{C_0 C_2^{\alpha}}.  \]
Hence $\mathcal{H}^{\alpha}_{\delta}(K) \ge  \frac{1}{C_0 C_2^{\alpha}}$ for every $\delta \in (0,\delta_2 /(C_2 + 1))$ and we obtain that $\mathcal{H}^{\alpha}(K) \ge  \frac{1}{C_0 C_2^{\alpha}}$. 

We see that $\mathcal{H}^{\alpha}(K) = \mathcal{H}^{\alpha}_{\infty}(K)  < \infty$ in the same manner as in the proof of Lemma \ref{lem:open-Haus-basic} below. 
Furthermore, we also see that $\overline{\dim}_M K \le \alpha$ by following the proof of \cite[Theorem 2.2.2 (ii)]{BishopPeres2017}. 
\end{proof}


\begin{Rem}
Conversely, if  Assumptions \ref{ass:two-additional-metric} and \ref{ass:additional-f}  hold and $0 < \mathcal{H}^{\beta}(K) < \infty$ for some $\beta > 0$, 
then $\sum_{j} r_j^{\beta} \ge 1$. 
If, in addition, $\mathcal{H}^{\beta}(f_k (K) \cap f_{\ell}(K)) = 0$ for all $k \ne \ell$, then $\sum_{j} r_j^{\beta} = 1$. 
See \cite[Theorem 3]{BessenyeiPenzes2022}\footnote{In the second statement of \cite[Theorem 3]{BessenyeiPenzes2022}, the assumption that each $f_j$ is surjective is missing.}. 
\end{Rem}



\section{The open set condition is necessary}\label{sec:osc}

In this section, we extend the necessity result proved in \cite{Schief1994} from the Euclidean setting to the setting of complete, normal and regular semimetric spaces.  

\begin{Thm}\label{thm:osc-necessary}
Suppose that Assumptions \ref{ass:two-additional-metric} and \ref{ass:additional-f} (ii) hold. 
Let $K$ be the attractor of $\{f_j\}_j$. 
Let $\alpha$ be the similarity dimension of the similitudes $\{f_j\}_j$. 
If $\mathcal{H}^{\alpha}(K) > 0$, then the open set condition holds.  
\end{Thm}


As an outline, we follow the proof in \cite[Section 9.6]{BishopPeres2017}, which is a simplification of  \cite{Schief1994}. 
We make only a few departures, but several parts are a little more complicated. 


We remark that the diameter of the $\epsilon$-neighborhood of $A$ is hard to estimate in terms of the diameter of $A$. 

For $A  \subset X$, $\beta \ge 0$ and $\delta \in (0, \infty]$, let 
\[ \mathcal{H}^{\beta}_{o;\delta}(A) \coloneqq \inf\left\{\sum_i \textup{diam}(U_i)^{\beta}  \colon A \subset \bigcup_i U_i, \ U_i \textup{ open}, \ \textup{diam}(U_i) \le \delta \right\}  \] 
and 
\[ \mathcal{H}^{\beta}_{o}(A) \coloneqq \lim_{\delta \to +0} \mathcal{H}^{\beta}_{o;\delta}(A). \] 
Then $\mathcal{H}^{\beta}_{o;\delta}(A) \ge \mathcal{H}^{\beta}_{\delta}(A)$ and hence $\mathcal{H}^{\beta}_{o}(A) \ge \mathcal{H}^{\beta}(A)$. 
$\mathcal{H}^{\beta}_{o}$ is an outer measure on $X$. 
By the same arguments as in the proof of  \cite[Lemma 11]{BessenyeiPenzes2022}, all Borel subsets of $X$ are $\mathcal{H}^{\beta}_{o}$-measurable. 

We say that two finite strings $\sigma, \tau \in \{1,\dots, \ell\}^{*}$ are {\it incomparable} if neither is a prefix of the other.  
The following corresponds to \cite[Proposition 2.1.3]{BishopPeres2017}. 

\begin{Lem}\label{lem:open-Haus-basic}
Let $\alpha$ be a positive real number such that $\sum_j r_j^{\alpha} = 1$. 
Then,\\
(i) $\mathcal{H}^{\alpha}_{o;\infty}(K) = \mathcal{H}^{\alpha}_{o}(K) < \infty$. \\
(ii) For every $ \mathcal{H}^{\alpha}_{o}$-measurable subset $B \subset K$, $\mathcal{H}^{\alpha}_{o;\infty}(B) = \mathcal{H}^{\alpha}_{o}(B)$. \\
(iii) If $\sigma$ and $\tau$ are incomparable, then $\mathcal{H}^{\alpha}_{o} \left(f_{\sigma}(K) \cap f_{\tau}(K) \right) = 0$. 
\end{Lem}

\begin{thebibliography}{99}
\bibitem[BaloghRohner2007]{BaloghRohner2007} Zolt\'an~M. Balogh and Heiner Rohner. \newblock Self-similar sets in doubling spaces. \newblock {\em Illinois J. Math.}, 51(4):1275--1297, 2007.
\bibitem[BessenyeiPales2017]{BessenyeiPales2017} Mih\'aly Bessenyei and Zsolt P\'ales. \newblock A contraction principle in semimetric spaces. \newblock {\em J. Nonlinear Convex Anal.}, 18(3):515--524, 2017.
\bibitem[BessenyeiPenzes2022]{BessenyeiPenzes2022} Mih{\'a}ly Bessenyei and Evelin P{\'e}nzes. \newblock Generalized fractals in semimetric spaces. \newblock {\em Nonlinear Anal., Theory Methods Appl., Ser. A, Theory Methods}, 220:15, 2022. \newblock Id/No 112853.
\bibitem[BishopPeres2017]{BishopPeres2017} Christopher~J. Bishop and Yuval Peres. \newblock {\em Fractals in probability and analysis}, volume 162 of {\em Cambridge Studies in Advanced Mathematics}. \newblock Cambridge University Press, Cambridge, 2017.
\bibitem[Bourbaki2007]{Bourbaki2007} Nicolas Bourbaki. \newblock {\em {\'E}l{\'e}ments de math{\'e}matique. {Topologie} g{\'e}n{\'e}rale. {Chapitres} 5 {\`a} 10}. \newblock Berlin: Springer, reprint of the 1974 original edition, 2007.
\bibitem[CJT2018]{CJT2018} Katarzyna Chrz{\k{a}}szcz, Jacek Jachymski, and Filip Turobo{\'s}. \newblock On characterizations and topology of regular semimetric spaces. \newblock {\em Publ. Math. Debr.}, 93(1-2):87--105, 2018.
\bibitem[KocsisPales2022]{KocsisPales2022} M\'aty\'as Kocsis and Zsolt P\'ales. \newblock Hutchinson's theorem in semimetric spaces. \newblock {\em J. Fixed Point Theory Appl.}, 24(4):Paper No. 75, 19, 2022.
\bibitem[Matkowski1975]{Matkowski1975} Janusz Matkowski. \newblock Integrable solutions of functional equations. \newblock {\em Diss. Math.}, 127, 1975.
\bibitem[Rus2001]{Rus2001} Ioan~A. Rus. \newblock {\em Generalized contractions and applications}. \newblock Cluj University Press, Cluj-Napoca, 2001.
\bibitem[Schief1994]{Schief1994} Andreas Schief. \newblock Separation properties for self-similar sets. \newblock {\em Proc. Am. Math. Soc.}, 122(1):111--115, 1994.
\end{thebibliography}
\end{document}
